\documentclass[10pt,twocolumn,a4paper]{ctexart}
\usepackage[margin=1.7cm,columnsep=0.65cm]{geometry}
\usepackage{amsmath,amssymb,bm}
\usepackage{booktabs}
\usepackage{graphicx}
\usepackage{microtype}
\usepackage{natbib}
\usepackage{xcolor}
\usepackage{hyperref}
\usepackage{enumitem}
\usepackage{fancyhdr}

\graphicspath{{../../artifacts/figures/}}
\setlength{\parindent}{2em}
\setlist{nosep,leftmargin=*}
\hypersetup{colorlinks=true,linkcolor=blue!55!black,
  citecolor=green!45!black,urlcolor=blue!65!black}
\pagestyle{fancy}
\fancyhf{}
\fancyhead[L]{Offline RL Interview Guide}
\fancyhead[R]{ShaneLiu04}
\fancyfoot[C]{\thepage}

\providecommand{\BestZeta}{0.9}
\providecommand{\IqlBestMediumMean}{--}
\providecommand{\IqlBestStitchMean}{--}
\providecommand{\AwacUmazeMean}{--}\providecommand{\AwacMediumMean}{--}
\providecommand{\IqlZAJUmazeMean}{--}\providecommand{\IqlZAJMediumMean}{--}
\providecommand{\IqlZACUmazeMean}{--}
\providecommand{\IqlZAJStitchMean}{--}\providecommand{\FbcBAStitchMean}{--}
\InputIfFileExists{../results_macros.tex}{}{}

\title{\vspace{-1.2em}\textbf{Offline Reinforcement Learning 工程与研究面试指南：\\
从分布偏移到 AWAC、IQL 与拼接}}
\author{ShaneLiu04\\
\small\url{https://github.com/ShaneLiu04}}
\date{}

\begin{document}
\maketitle
\thispagestyle{fancy}

\begin{abstract}
本文整理 offline RL 岗位面试里反复出现的一条主线：固定数据集上的 Bellman
备份为什么会外推，各类算法如何把策略改进限制在数据支撑内，以及这些选择在
工程上对应哪些超参、日志和失败模式。内容覆盖 batch RL、行为克隆、BCQ、BEAR、
CQL、AWR/AWAC、CRR、IQL、expectile、模型类离线方法和序列建模基线，并用 AntMaze
与 PointMass 上的实测数字说明每个机制改变的是哪一项误差。本文面向算法工程与
研究岗位：既要能写出损失，也要能判断一份日志里的高 Q 值是学到了拼接还是
过估计。
\end{abstract}

\textbf{关键词：}Offline RL；Distribution Shift；AWAC；IQL；CQL；Stitching；Expectile

\section{固定数据集改变了什么}
Online RL 的策略改进依赖新的交互：价值估高的动作会被执行，环境返回真实的
后继，错误的外推会被样本纠正\citep{sutton2018rl}。Offline RL 去掉了这一回路。
数据集 $\mathcal{D}$ 在训练开始前就固定，算法只能在这些 $(s,a,r,s')$ 上更新
\citep{lange2012batch,levine2020offline}。工业场景里这是常态：日志来自人工、
脚本或旧策略，在线试错的代价过高，例如机器人损坏、推荐系统的用户流失、
自动驾驶的安全约束。

行为克隆把问题收成监督学习，$\max_\psi\mathbb{E}_{\mathcal{D}}[\log\pi_\psi(a\mid s)]$
\citep{pomerleau1991bc}。它没有外推，也没有改进。部署时状态分布一旦偏离采集分布，
误差会沿时间复合，这就是 covariate shift\citep{ross2011dagger}。Offline RL 想在
不采集新数据的前提下，得到比行为策略更高的回报。它必须利用动态规划把不同轨迹的
片段接起来，同时避免在没有数据的动作上相信自己的 Q 值。

\section{外推误差}
标准 Q-learning 的目标含有 $\max_{a'}Q(s',a')$ 或 $\mathbb{E}_{a'\sim\pi}[Q(s',a')]$。
函数近似在数据稀疏的区域没有约束，最大化或策略采样会主动挑选这些高估点。
下一次备份把高估写回有数据的状态，误差沿轨迹反向传播。Fujimoto 等人把这种现象
称为 extrapolation error，并指出 off-policy deep RL 在完全没有探索时会发散
\citep{fujimoto2019bcq}。Kumar 等人进一步把备份误差写成 bootstrapping error：
不可见动作的价值误差被 $\gamma$ 折扣后送回可见状态，策略若增加这些动作的概率，
误差会被放大\citep{kumar2019bear}。

因此 offline 算法都在回答同一个设计问题：策略改进的提议分布，和价值备份的
查询分布，哪一个允许离开 $\mathcal{D}$。只约束其中一个通常不够。TD3 的 clipped
double-Q 减轻在线设置里的过估计\citep{fujimoto2018td3,vanhasselt2016ddqn}，但两个
网络可以在同一个未见动作上一起高估，所以它不是 offline 约束。

\section{约束与悲观}
第一类方法显式限制策略接近行为分布。BCQ 用一个条件 VAE 拟合行为策略，动作提案
只在 VAE 样本的小扰动内做最大化\citep{fujimoto2019bcq}。BEAR 用 MMD 把学习策略
撑在行为策略的支撑上，而不是逐动作匹配密度，从而允许多模态数据
\citep{kumar2019bear}。BRAC 直接在策略目标里加行为策略的散度惩罚
\citep{wu2019behavior}。TD3+BC 把这个思想收成一项极简正则：在 TD3 的策略损失上
加 $\lambda_{\mathrm{BC}}\|a-\pi(s)\|^{2}$，并用 Q 的尺度做归一，避免正则权重随
奖励尺度漂移\citep{fujimoto2021minimalist}。面试里要能说明：支撑约束保留拼接
所需的多模态，密度匹配会把策略拉向次优模态的平均。

第二类方法不改提案分布，而把未见动作的价值压低。CQL 在 Q 损失里加入
$\mathbb{E}_{a\sim\mu}[Q(s,a)]-\mathbb{E}_{a\sim\mathcal{D}}[Q(s,a)]$ 一类的正则，
使任意策略下的 Q 值下界尽量不超过数据动作上的 Q\citep{kumar2020cql}。它的
工程代价是需要对 $\mu$ 采样，正则系数对任务敏感。模型类方法把不确定区域的奖励
改小：MOPO 用动力学集成的方差惩罚奖励\citep{yu2020mopo}，MOReL 把不确定转移
送进吸收的低价值状态\citep{kidambi2020morel}，COMBO 把 CQL 式保守和模型 rollout
合在一起\citep{yu2021combo}。模型误差本身会变成新的外推源，所以惩罚必须和
集成分歧校准，而不是手写一个常数\citep{ghasemipour2022why}。

\section{加权回归：AWR、AWAC 与 CRR}
Advantage-weighted regression 把策略改进写成带权的极大似然
\citep{peng2019awr}。从 KL 约束的策略改进出发，最优解满足
\begin{equation}
\pi_{k+1}(a\mid s)\propto \pi_{\beta}(a\mid s)\exp\bigl(A^{\pi_k}(s,a)/\lambda\bigr).
\end{equation}
当行为数据就是 $\pi_{\beta}$ 的样本时，投影回参数族等价于最小化
\begin{equation}
-\mathbb{E}_{(s,a)\sim\mathcal{D}}\bigl[\log\pi_\psi(a\mid s)\exp(A/\lambda)\bigr].
\end{equation}
$\lambda$ 是温度，也是离开行为分布的步长。它小，则少数高优势样本主导梯度；
它大，则退化为行为克隆。权重必须 detach，否则优势会为了减小损失而被网络改写。
实现里还要把 $\exp(A/\lambda)$ 截断，避免单个离群优势把梯度打爆。

AWAC 用同一项损失，并在离线数据上先训练，再在线微调\citep{nair2020awac}。
价值备份仍是 $a'\sim\pi(\cdot\mid s')$ 的 expected SARSA，配合双 Q 与目标网络。
因此 AWAC 的策略步是保守的，价值步并不保守。这解释了它在短迷宫上可以工作、
在长时域稀疏奖励上容易被早期随机 $a'$ 带偏的现象。CRR 做了更明确的过滤：
权重可以是 $\exp(A/\beta)$，也可以是只保留正优势的指示函数
\citep{wang2020crr}。面试时把这三项放在一起：AWR 给出闭式动机，AWAC 给出
离线到在线的训练日程，CRR 给出“是否丢弃负优势样本”的消融维度。

本文的 AWAC 在 u-maze 上最终回报为 $\AwacUmazeMean$，在 medium-diverse 上为
$\AwacMediumMean$。温度扫描用来检查默认 $\lambda=0.1$ 是否落在模仿和改进之间。

\section{IQL：不查询未见动作}
IQL 的目标是让价值备份完全 in-sample\citep{kostrikov2022iql}。它引入 $V(s)$，
用 expectile 回归拟合目标 Q 的上尾，而不是均值\citep{newey1987expectile}。
损失 $L_2^{\zeta}(\mu)=\lvert\zeta-\mathbf{1}\{\mu<0\}\rvert\mu^{2}$ 在 $\zeta>0.5$ 时
给正残差更大的权重。一阶条件是
\begin{equation}
\zeta\,\mathbb{E}[(Q-V)_{+}]=(1-\zeta)\,\mathbb{E}[(V-Q)_{+}].
\end{equation}
$\zeta=0.5$ 时 $V$ 是条件均值；$\zeta\to 1$ 时 $V$ 逼近数据中该状态上的最大 Q。
Q 的目标改成 $r+\gamma V(s')$，并且 $V$ 不通过这条损失接收 Q 的梯度。于是 Q 网络
从未在策略采样的 $a'$ 上被训练。策略抽取再使用上面的指数加权回归，优势为
$Q-V$，温度的倒数就是论文里的 $\beta$。

这和 CQL 的差别值得在面试里讲清楚。CQL 主动压低未见动作的 Q，需要额外采样。
IQL 根本不给未见动作定义备份，逼近的是数据集条件分布的上 expectile。因此 IQL
能拼接的是“数据里真实出现过的高价值后继”，不能发现数据支撑之外的捷径。
$\zeta$ 太低会学到差轨迹的均值，太高会把偶然的高回报噪声当成价值。
本文 u-maze 上 $\zeta=0.2$ 的回报是 $\IqlZACUmazeMean$，$\zeta=0.9$ 是
$\IqlZAJUmazeMean$。medium 上选用 $\zeta=\BestZeta$，回报 $\IqlBestMediumMean$。

后续工作把这一思想再参数化。XQL 用极值理论连接 max-entropy 与 expectile，
给出 Gumbel 误差下的损失\citep{garg2023xql}。IDQL 保留 IQL 的价值学习，把
高斯策略换成 diffusion，以覆盖行为分布的多模态\citep{hansenestruch2023idql}。
Fisher-BRC 则从散度正则的角度约束 critic，而不是换备份算子
\citep{kostrikov2021fbrc}。DR3 指出深度价值网络即使在 in-sample 回归下也会
因为隐式正则把 Q 推向共同的低秩方向，需要显式惩罚特征内积
\citep{kumar2022dr3}。这些论文回答的是“IQL 已经不查询 OOD 动作之后，还剩什么误差”。

\section{拼接与监督学习基线}
Stitching 指数据集中不存在最优轨迹，但最优轨迹的每一段都出现在某条次优轨迹里。
动态规划能把这些段接上，因为 $V(s)$ 总结的是后继的价值，而不是“必须沿同一条
轨迹走下去”。行为克隆和 Filtered BC 没有这个总结，它们只提高所见动作的似然。
只保留最高回报的 $10\%$ 轨迹可以去掉明显的失败样本，但保留下来的轨迹本身仍是
次优的，策略无法走出它们的凸组合之外。

Return-conditioned supervised learning 用目标回报当条件，直接拟合
$p(a\mid s,R)$\citep{emmons2022rvs,ghosh2021gcsl}。Decision Transformer 把整段
轨迹当成序列，用自回归模型按目标 return-to-go 生成动作\citep{chen2021dt}；
Trajectory Transformer 在同一框架里离散化状态和动作\citep{janner2021tt}。
这类方法实现简单，在有完整成功轨迹的数据上很强。它们弱在拼接：模型看到的训练
序列并没有把 A 轨迹的前半和 B 轨迹的后半接在同一个 context 里，除非数据里已经
有这样的拼接样本。AntMaze 因此成为区分动态规划和序列克隆的常用探针
\citep{fu2020d4rl}。

本文的 PointMass 数据集最大回报为 $-46$、平均约 $-104$。选用 $\zeta$ 的 IQL
平均评估回报为 $\IqlBestStitchMean$，Filtered BC 为 $\FbcBAStitchMean$。比较对象应当是 $-46$，
而不是 $0$：超过数据集最大值才说明出现了数据中不存在的更好轨迹。

\section{工程清单}
面试里的实现题通常不考新损失，而考下面这些容易写错的接口。

\begin{itemize}
  \item 奖励约定要和评估一致。D4RL AntMaze 原始目标奖励常为 $1$，训练代码把它减
        $1$ 后目标步为 $0$。评估若忘记同一平移，曲线会整体偏移一个回合长度。
  \item 终止标记只在真正结束时截断 bootstrap。超时（time limit）不应当成
        $d=1$，否则长回合的价值被系统性压低。v0 与 v2 迷宫数据的一个已知差别
        就是早期终止的处理\citep{fu2020d4rl}。
  \item 目标网络用 EMA，备份进 $\mathrm{no\_grad}$。IQL 里 $V(s')$ 也要断开，
        否则 Q 和 V 会互相追逐。
  \item 指数权重先 detach 再截断。优势的标准化可以再降方差，但会改变 $\lambda$
        的语义，报告里必须写明有没有做。
  \item 连续动作在送进 Q 之前截断到环境范围。高斯尾部的样本会让 Q 在
        $\lvert a\rvert>1$ 的区域接收梯度。
  \item 三种子的均值不要和单次评估的轨迹标准差混用。跨种子统计才对应算法稳定性
        \citep{agarwal2021precipice}。
  \item 数据集文件要校验 HDF5 超块里的 \texttt{stored\_eof}。镜像上的
        medium-diverse v0 文件声明长度为 $231\,095\,454$ 字节，实体只有
        $100\,917\,248$ 字节。这种文件能通过下载大小检查，但不能被 h5py 打开。
\end{itemize}

日志上至少保留评估回报、成功率、Q 损失、V 损失和优势权重的分位数。Q 损失下降
而成功率不升，通常是价值在拟合一个不会到达目标的行为分布。优势权重的最大值
长期顶在截断上限，说明 $\lambda$ 过小或优势尺度没有控制。

从离线到在线的微调是另一类常见岗位问题。AWAC 的原始动机就是用离线数据初始化，
再用少量在线样本改进\citep{nair2020awac}。Cal-QL 在离线 CQL 预训练和在线微调
之间校准 Q 的尺度，避免上线后价值崩溃\citep{nakamoto2023calql}。RLPD 则把离线
数据和高 UTD 的在线更新放进同一个 buffer，强调对称采样而不是复杂的悲观正则
\citep{ball2023rlpd}。讲这条线时要区分：离线阶段防的是外推，上线阶段防的是
分布突然变了之后旧 Q 的尺度不再匹配。

\section{怎么读一份实验}
D4RL 把不同任务的回报映射到归一化分数，AntMaze 的参考上界对应成功率
\citep{fu2020d4rl}。RL Unplugged 使用另一套数据划分，数字不能直接和 D4RL 对比
\citep{gulcehre2020rlu}。CORL 提供了研究向的统一实现，适合作为复现时的对照
代码，而不是作为黑盒替换自己的备份算子\citep{tarasov2022corl}。

读本文的 AntMaze 曲线时，先看最终三种子均值，再看成功率是否和回报同向。
回报上升但成功率不变，可能只是失败轨迹稍短。读 PointMass 时，先把结果和
数据集最大值 $-46$ 放在一起，再看轨迹图里评估路径是否穿过了单条数据轨迹没有
完整覆盖的走廊。expectile 和 $\lambda$ 的扫描应呈现单峰或平台：两端变差说明旋钮
有效，单调则说明搜索范围没有覆盖到退化和过拟合的两侧。

综述可以用来组织阅读顺序，而不能代替原始损失的推导\citep{prudencio2023survey}。
一个够用的阅读顺序是：extrapolation error\citep{fujimoto2019bcq}，支撑约束
\citep{kumar2019bear}，保守 Q\citep{kumar2020cql}，加权回归
\citep{peng2019awr,nair2020awac,wang2020crr}，in-sample expectile
\citep{kostrikov2022iql}，然后才是 diffusion 策略、序列模型和在线微调。

\section{岗位问题上的回答方式}
“为什么不能直接把 DQN 或 SAC 套到日志上？”因为最大化或熵正则策略会查询
$\mathcal{D}$ 没有覆盖的动作，函数近似在那里没有惩罚
\citep{mnih2015dqn,haarnoja2018sac,fujimoto2019bcq}。

“AWAC 和 IQL 的策略损失一样，为什么结果不同？”因为差异在 critic。AWAC 的
$y$ 含 $a'\sim\pi$，IQL 的 $y$ 只含 $V(s')$。长时域上这个差别会被 $\gamma^{T}$
放大。

“$\zeta$ 和 $\lambda$ 分别调什么？”$\zeta$ 选择价值去贴数据上尾的程度，$\lambda$
选择策略偏离行为样本的程度。两个都变大并不等价：高 $\zeta$ 加低 $\lambda$ 会在
一个很尖的上尾上做很大的策略步，方差最大。

“如何证明出现了拼接？”需要一个已知最优轨迹不在数据集中的环境，报告数据最大
回报和评估回报，并给出轨迹可视化。只有平均回报提高、没有和数据上限比较，
不能区分“模仿了最好的那几条”和“拼出了更好的一条”。

\bibliographystyle{plainnat}
\bibliography{../references}

\end{document}
